% Source: https://artofproblemsolving.com/community/c3h1378409_paper_problem
% Stable identity: aops-topic-1378409-q01
Put the centre at $(0,0)$, so the paper is $-11\le x\le11$, $-10\le y\le10$. The segments have lengths $0.5,1.5,2.5,\dots$ and directions north, east, south, west repeating. After the first four segments the pencil is at
\[
(0,0)\to(0,0.5)\to(1.5,0.5)\to(1.5,-2)\to(-2,-2),
\]
and in general each group of four moves the pencil by $(-2,-2)$: after $4k$ segments it is at $(-2k,-2k)$. Inside the group starting there, the segments of lengths $4k+0.5$, $4k+1.5$, $4k+2.5$, $4k+3.5$ end at
\[
(-2k,\,2k+0.5),\quad(2k+1.5,\,2k+0.5),\quad(2k+1.5,\,-2k-2),\quad(-2k-2,\,-2k-2).
\]
The largest coordinate reached so far grows slowly, so we look for the first segment that would leave the sheet. For $k=4$ the group ends at $(-10,-10)$ with all points inside ($2k+0.5=8.5\le10$, $2k+1.5=9.5\le11$, $2k+2=10\le10$ and $\le11$). For $k=5$ the first segment starts at $(-10,-10)$ and goes north with length $20.5$; it would end at $y=10.5>10$, so it reaches the north edge after $20$ inches, and Bob stops there.

So he draws the $20$ complete segments of lengths $0.5,1.5,\dots,19.5$ and then $20$ inches of the next one:
\[
(0.5+1.5+\cdots+19.5)+20=\dfrac{20\,(0.5+19.5)}{2}+20=200+20=\boxed{220}\text{ inches}.
\]
